% !TEX root = ../main.tex
\section{Q04: De 5 incitamentkonflikter -- principal-agent under costly contracting og asymmetrisk informasjon}\label{q:04}

\begin{quote}
\emph{Med udgangspunkt i pensums definition af en virksomhed bedes du redegøre for 5 grundlæggende incitamentkonflikter mellem ejere og virksomhedsledere. Herefter bedes der redegjort for, hvordan incitamentkonflikter mellem en principal og en agent (antag to virksomheder) kan organiseres og styres ved hjælp af kontrakter under forudsætning af: costly contracting og asymmetrisk information. Begreberne præ- og post-kontraktuelle problemer (adverse selection og moral hazard) bedes inddraget i denne redegørelse.}
\end{quote}

\subsection*{Svar (kort)}
De fem incitamentkonfliktene mellom eier og leder -- shirking, perks/lønn, risikoaversjon, korttidshorisont og overinvestering -- er de samme som i Q03. Det som skiller Q04 er forutsetningene: når kontrakter er kostbare å skrive/overvåke/håndheve og informasjonen er asymmetrisk, kan kontrakten ikke lenger eliminere konfliktene. Vi havner i en \emph{second-best}-løsning der pre-kontraktuelle problemer (adverse selection -- skjult \emph{type}) håndteres med screening og signaling, mens post-kontraktuelle problemer (moral hazard -- skjult \emph{handling}) håndteres med monitoring, bonding og insentivlønn -- og en residuelt loss aksepteres. Når prinsipal og agent er to virksomheter (interfirm contracting) løftes problemet til kontrakter om leveranser, partnerskap og felles investeringer. [SLIDE]

\subsection*{Relevante teorier (kort)}

\paragraph{\thref{incentive-conflicts-owner-manager}{De fem incitamentkonfliktene}\quad\SThree}
Samme liste som i Q03: (1) shirking (innsatsreduksjon), (2) salaries and perks (overforbruk av personalegoder), (3) risikoaversjon (leder unnviker NPV-positive risikable prosjekter), (4) ulik tidshorisont (kortsiktig profitfokus), (5) overinvestering / empire building. Konfliktene er reelle nettopp fordi det kontraktuelle apparatet ikke er gratis -- det er forutsetningen om \emph{costly contracting} som gjør dem til organisasjonsproblemer. [SLIDE -- BSZ kap. 10]

\paragraph{\thref{principal-agent}{Principal-Agent-teori}\quad\SThree}
Modellen $W = W_0 + \beta Q$ med $Q = \alpha e + \mu$ blir nå en \emph{second-best}-modell: optimal $\beta$ avveier motivasjon mot risikopremie og kontraktomkostninger. Når prinsipal og agent er to virksomheter, er $\alpha$ partnerens reelle kompetanse (skjult før kontrakt), $e$ er innsatsen partneren legger inn (skjult etter kontrakt), og $\mu$ er felles markedsstøy. Optimal kontrakt blir aldri first-best. [SLIDE]

\paragraph{\thref{costly-contracting}{Costly Contracting}\quad\SOne/\SThree}
Search-, forhandlings-, drafting-, monitoring- og enforcement-omkostninger gjør komplette kontrakter umulige. Kontrakter blir \emph{ufullstendige}, og residuelle beslutningsrettigheter må allokeres -- enten via eierskap (vertikal integrasjon) eller via reputasjon og implisitte avtaler. Forutsetningen er kontrasten mot Q03's \thref{costless-contracting}{costless contracting}, der konfliktene kunne skrives bort i en first-best-kontrakt. [SLIDE/BOK]

\paragraph{\thref{adverse-selection}{Adverse Selection (pre-kontraktuelt -- hidden information)}\quad\SThree}
Før kontrakten signeres kjenner én part sin egen \emph{type} (kvalitet, kompetanse, risiko) bedre enn motparten. I to-virksomhetsperspektivet: en teknologileverandør vet om elektrolyseteknologien faktisk er moden. Akerlofs lemons-mekanisme får markedet til å selektere de dårlige inn. Løsninger: \emph{signaling} (garantier, sertifiseringer, eierandel), \emph{screening} (due-diligence, kontraktsmenyer som trigger selvseleksjon), reputasjon. [SLIDE]

\paragraph{\thref{moral-hazard}{Moral Hazard (post-kontraktuelt -- hidden action)}\quad\SThree}
Etter kontrakten kan agenten ta skjulte handlinger -- shirke, kutte kvalitet, fortie problemer, prioritere egen profit over felles. Løsninger: \emph{monitoring} (revisjon, KPI-oppfølging, rapportering), \emph{bonding} (skin-in-the-game, vesting, claw-back), variabel lønn ($\beta > 0$). En \emph{residual loss} gjenstår alltid. Summen utgjør \thref{agency-costs}{agency costs}. [SLIDE]

\subsection*{Illustrasjoner fra forelesning}

\begin{figure}[h]
  \centering
  \includegraphics[width=0.5\linewidth]{BSZ_kap_8_2026/by_slide/slide_003/BSZ_kap_8_2026_slide_003_asset_01_image2.jpeg}
  \caption{Transaksjonskostnader: søk-, forhandlings- og håndhevingskostnader som gjør kontrakter ufullstendige -- grunnlaget for costly contracting i Q04. (BSZ kap.\ 8, slide 3) [SLIDE]}
\end{figure}

\begin{figure}[h]
  \centering
  
\includegraphics[width=0.7\linewidth]{BSZ_kap_10_2026/by_slide/slide_005/BSZ_kap_10_2026_slide_005_asset_01_image2.png}
  \caption{Virksomheten som en \emph{nexus of contracts} -- kontrakter mellom eiere, ledere, ansatte, leverandører, kunder, banker. Q04 fokuserer på principal-agent-kontrakten mellom to virksomheter under costly contracting. [SLIDE -- BSZ kap. 10, slide 5]}
\end{figure}

\begin{center}
\renewcommand{\arraystretch}{1.15}
\begin{tabular}{@{}p{3.6cm}p{3.6cm}p{6.5cm}@{}}
\toprule
\textbf{Problem} & \textbf{Tidspunkt} & \textbf{Hovedløsning} \\
\midrule
Adverse selection & Før kontrakt (skjult type) & Signaling, screening, due-diligence, reputasjon \\
Moral hazard & Etter kontrakt (skjult handling) & Monitoring, bonding, insentivkontrakt ($\beta$) \\
Residual loss & Alltid & Aksepteres -- second-best \\
\bottomrule
\end{tabular}
\end{center}
\textit{Tabell: Pre- vs.\ post-kontraktuelle problemer og hvordan de håndteres. [SLIDE]}

\subsection*{Kontrakter under costly contracting og asymmetrisk informasjon -- kort}
\begin{itemize}\setlength{\itemsep}{2pt}
\item \textbf{Pre-kontraktuelt (adverse selection).} Kjøperen (prinsipal) kan ikke skille gode og dårlige tilbydere. Verktøy: due-diligence, sertifisering, garantier, kontraktsmenyer (screening), reputasjon, prøvebestillinger. [SLIDE]
\item \textbf{Post-kontraktuelt (moral hazard).} Etter signering kan agenten redusere kvalitet/innsats skjult. Verktøy: milepælsbetaling, performance bonds, monitoring/revisjon, insentivkontrakter ($W = W_0 + \beta Q$), claw-back, vertikal integrasjon ved særlig hold-up-risiko. [SLIDE]
\item \textbf{Residual control.} Når kontrakten er ufullstendig (begrenset rasjonalitet), er det den som eier aktiva som bestemmer i uregulerte tilstander -- begrunnelsen for vertikal integrasjon. [BOK]
\item \textbf{Trade-off:} Optimal $\beta$ balanserer motivasjon mot risikopremie pluss monitoring/bonding-kostnader. Resultatet er second-best, ikke first-best. [SLIDE]
\end{itemize}

\subsection*{Fordele og ulemper ved kontraktsmekanismene}
\begin{itemize}\setlength{\itemsep}{2pt}
\item \textbf{Due-diligence / screening.} Fordel: avdekker skjult type før kontrakt; reduserer adverse selection. Ulempe: kostbart, gir bare punktestimat, kan ikke fange dynamiske endringer; risiko for ``vinnerens forbannelse'' der den som vinner budet er den som overvurderer mest. [SLIDE]
\item \textbf{Signaling (garantier, sertifisering, eierandel).} Fordel: informert part forplikter seg på troverdig vis (separating equilibrium). Ulempe: signalet er kostbart for høye-typer også; svake typer kan etterligne hvis kostnaden ikke er separerende nok. [SLIDE]
\item \textbf{Monitoring.} Fordel: gjør skjult handling observerbar etter kontrakt; reduserer moral hazard. Ulempe: direkte kostnad pluss demotiverende mikroledelse; måler ofte input, ikke output, og trigger gaming. [SLIDE]
\item \textbf{Bonding (skin-in-the-game, vesting, claw-back).} Fordel: agenten internaliserer prinsipalens risiko -- selvbinding er billigere enn ekstern overvåkning. Ulempe: krever at agenten har formue å binde, og overfører risiko til typisk risikoavers part (krever risikopremie). [BOK]
\item \textbf{Insentivkontrakt ($W = W_0 + \beta Q$).} Fordel: lineær mekanisme som er enkel å forstå og håndheve. Ulempe: når $Q$ inneholder mye støy, blir nødvendig $\beta$ liten -- svake insentiver er det beste man kan oppnå. [SLIDE]
\item \textbf{Vertikal integrasjon / residual control.} Fordel: eierskap allokerer beslutningsrett i uregulerte tilstander; løser hold-up når aktivaspesifisitet er høy. Ulempe: bytter markedsdisiplin mot intern byråkrati og svake insentiver i den integrerte enheten. [BOK -- BSZ kap. 19]
\end{itemize}

\subsection*{Forutsetninger for valg av mekanisme}
\begin{itemize}\setlength{\itemsep}{2pt}
\item \textbf{Adverse selection dominerer} (typeusikkerhet er hovedproblemet): screening og signaling først; gjentatte transaksjoner gir reputasjonsverdi.
\item \textbf{Moral hazard dominerer} (handlingen er hovedproblemet): monitoring + variabel $\beta$ + bonding; jo mer presis $Q$-måling, desto høyere $\beta$.
\item \textbf{Begge til stede} (typisk virkelighet): kombiner screening (forkant) + insentivkontrakt (etterkant); designvalget styres av relative kostnader.
\item \textbf{Aktivaspesifisitet er høy + kontrakt er ufullstendig:} integrer vertikalt; bytt eksterne insentiver mot interne residuelle rettigheter.
\item \textbf{Agenten er svært risikoavers eller fattig:} reduser $\beta$, øk $W_0$, fokuser på monitoring i stedet for bonding.
\end{itemize}

\subsection*{Real-world eksempler}

\textbf{Eks.~1 -- Statoil--Hydro-fusjonen (2007), adverse selection.}\quad [EKS]
To store norske virksomheter med privat informasjon om egne reserver, kontrakter og forpliktelser måtte gjennom flere måneder med due-diligence (klassisk screening) og uavhengige verifikasjoner (signaling) før bytteforholdet kunne fastsettes. Selv etter dette viste integrasjonen flere ``overraskelser'' -- typiske residual losses fra ikke fullstendig løst pre-kontraktuell informasjonsasymmetri.

\textbf{Eks.~2 -- DNB rådgiverprovisjoner (2010-tallet), moral hazard -- sammenligning.}\quad [EKS]
DNB-rådgivere fikk provisjon for å selge bankens egne fond. Atferden var skjult både for kunde og toppledelse, og rådgiverne valgte produkter med høyere kostnad enn nødvendig -- klassisk post-kontraktuell moral hazard. Forbrukerrådets shadow shopping (monitoring) avslørte forholdet og førte til at provisjonsmodellen ble fjernet (regulatorisk endring av $\beta$-strukturen). Sammenligning med Statoil-Hydro: samme principal-agent-rammeverk, men problemet er på \emph{ulike sider} av kontraktinngåelsen -- adverse selection vs.\ moral hazard. Løsningsverktøyene er derfor ulike: due-diligence og signaling i forkant vs.\ monitoring og endret insentivkontrakt i etterkant.

\subsection*{Kobling til Mærsk Power-to-X}
Mærsks Power-to-X-satsing er et lærebok-eksempel på begge problemene under costly contracting mellom \emph{to virksomheter}. \emph{Adverse selection}: ved valg av teknologipartner (elektrolyseleverandør, hydrogenkjøper) kjenner partneren sin egen reelle modenhet og kostnadsstruktur bedre enn Mærsk -- løses med teknisk due-diligence, krav om referanseprosjekter (signaling) og garantier. \emph{Moral hazard}: etter kontraktinngåelse kan en EPC-kontraktør skjule fremdriftsproblemer eller kvalitetskutt -- løses med milepælsbetaling, on-site supervision (monitoring), performance bonds (bonding) og deling av oppside via partnerskap (variabel $\beta$). Når kontrakten ikke kan dekke alle fremtidige tilstander (teknologi-, regulatorisk- og prisrisiko), kan Mærsk velge vertikal integrasjon på de mest spesifikke aktivaene. [CASE]

\subsection*{Fallgruber}
\begin{itemize}\setlength{\itemsep}{2pt}
\item Ikke gjenta Q03 -- her er hele poenget at \emph{forutsetningene} er andre (costly + asymmetrisk).
\item Forveksle adverse selection og moral hazard -- nøkkelen er \emph{tidspunkt} og \emph{type skjult variabel} (egenskap vs.\ handling).
\item Glemme residual loss -- selv optimal kontrakt etterlater et tap.
\item Hoppe over to-virksomhets-rammen -- spørsmålet ber eksplisitt om interfirm-kontekst.
\item Gi inntrykk av at kontrakten kan ``løse'' problemet -- under costly contracting reduseres det, ikke elimineres.
\item Glemme at signaling og screening er to forskjellige svar (informert part vs.\ uinformert part tar grep).
\end{itemize}

\subsection*{Håndskrivbar disposisjon (15-min forberedelse)}
\begin{itemize}\setlength{\itemsep}{2pt}
\item \textbf{Tese:} Samme 5 konflikter som Q03, men under costly contracting + asymmetrisk info kan kontrakten ikke eliminere dem -- kun second-best.
\item Virksomheten = \emph{nexus of contracts}; her ser vi på principal-agent-kontrakten mellom \emph{to virksomheter}.
\item \thref{incentive-conflicts-owner-manager}{5 konflikter}: shirking, perks, risikoaversjon, tidshorisont, overinvestering.
\item \thref{costly-contracting}{Costly contracting}: search/forhandling/drafting/monitoring/enforcement -- kontrakten blir ufullstendig.
\item \thref{adverse-selection}{Adverse selection} = pre-kontraktuell, skjult \emph{type} $\Rightarrow$ signaling (informert part) + screening / due-diligence (uinformert part).
\item \thref{moral-hazard}{Moral hazard} = post-kontraktuell, skjult \emph{handling} $\Rightarrow$ monitoring + bonding + insentivkontrakt ($\beta > 0$).
\item \thref{principal-agent}{Principal-agent}-modellen $W = W_0 + \beta Q$: optimal $\beta$ balanserer motivasjon mot risikopremie + agency cost.
\item Residual loss $+$ monitoring $+$ bonding = \thref{agency-costs}{agency costs}; transaksjonen kan falle vekk hvis disse blir for store.
\item Eks. 1: Statoil-Hydro 2007 -- adverse selection, due-diligence.
\item Eks. 2: DNB rådgiverprovisjoner -- moral hazard, monitoring + regulatorisk endring av $\beta$.
\item Mærsk Power-to-X: teknologipartner = adverse selection; EPC-kontraktør = moral hazard; ufullstendig kontrakt $\Rightarrow$ vertikal integrasjon for kjerneaktiva.
\item Kontrast Q03: der \thref{costless-contracting}{costless contracting} = first-best; her costly + asymmetrisk = second-best.
\item Søyle: \SThree\ primært, også \SOne\ (residual control rights).
\end{itemize}
